\begin{ex}
Tính: $25\text{ kg} + 14\text{ kg}$.
\shortans{39}
\loigiai{
Ta có: $25 + 14 = 39$ kg.
}
\end{ex}

\begin{ex}
Tính: $48\text{ kg} - 23\text{ kg}$.
\shortans{25}
\loigiai{
Ta có: $48 - 23 = 25$ kg.
}
\end{ex}

\begin{ex}
Tính: $32\text{ }l + 16\text{ }l$.
\shortans{48}
\loigiai{
Ta có: $32 + 16 = 48$ lít.
}
\end{ex}

\begin{ex}
Tính: $57\text{ }l - 35\text{ }l$.
\shortans{22}
\loigiai{
Ta có: $57 - 35 = 22$ lít.
}
\end{ex}

\begin{ex}
Tính giá trị biểu thức: $18\text{ kg} + 12\text{ kg} - 10\text{ kg}$.
\shortans{20}
\loigiai{
Ta có $18 + 12 = 30$, rồi $30 - 10 = 20$ kg.
}
\end{ex}

\begin{ex}
Tính giá trị biểu thức: $45\text{ }l - 20\text{ }l + 15\text{ }l$.
\shortans{40}
\loigiai{
Ta có $45 - 20 = 25$, rồi $25 + 15 = 40$ lít.
}
\end{ex}

\begin{ex}
Điền số thích hợp: $28\text{ kg} + \dots\text{ kg} = 49\text{ kg}$.
\shortans{21}
\loigiai{
Số cần điền là: $49 - 28 = 21$.
}
\end{ex}

\begin{ex}
Điền số thích hợp: $65\text{ }l - \dots\text{ }l = 30\text{ }l$.
\shortans{35}
\loigiai{
Số cần điền là: $65 - 30 = 35$.
}
\end{ex}

\begin{ex}
Điền số thích hợp: $\dots\text{ kg} - 14\text{ kg} = 35\text{ kg}$.
\shortans{49}
\loigiai{
Số cần điền là: $35 + 14 = 49$.
}
\end{ex}

\begin{ex}
Một con chó nặng 14 kg, con mèo nhẹ hơn con chó 10 kg. Con mèo nặng bao nhiêu ki-lô-gam?
\shortans{4}
\loigiai{
Con mèo nặng: $14 - 10 = 4$ kg.
}
\end{ex}

\begin{ex}
Cả con chó và con mèo nặng tất cả bao nhiêu ki-lô-gam?
\shortans{18}
\loigiai{
Cả hai con nặng: $14 + 4 = 18$ kg.
}
\end{ex}

\begin{ex}
Thùng thứ nhất có 38 lít dầu, thùng thứ hai có ít hơn thùng thứ nhất 15 lít dầu. Thùng thứ hai có bao nhiêu lít dầu?
\shortans{23}
\loigiai{
Thùng thứ hai có: $38 - 15 = 23$ lít.
}
\end{ex}

\begin{ex}
Cả hai thùng có tất cả bao nhiêu lít dầu?
\shortans{61}
\loigiai{
Cả hai thùng có: $38 + 23 = 61$ lít.
}
\end{ex}

\begin{ex}
Bao thóc thứ nhất nặng 45 kg, bao thóc thứ hai nặng nhiều hơn bao thứ nhất 12 kg. Bao thóc thứ hai nặng bao nhiêu ki-lô-gam?
\shortans{57}
\loigiai{
Bao thứ hai nặng: $45 + 12 = 57$ kg.
}
\end{ex}

\begin{ex}
Một can chứa đầy 20 lít nước. Sau khi dùng 8 lít, người ta lại đổ thêm vào can 5 lít nữa. Can nước lúc này có bao nhiêu lít?
\shortans{17}
\loigiai{
Còn lại: $20 - 8 + 5 = 17$ lít.
}
\end{ex}

\begin{ex}
Một quả dưa hấu nặng 4 kg, một nải chuối nặng 3 kg. Cả hai thứ nặng bao nhiêu ki-lô-gam?
\shortans{7}
\loigiai{
Cả hai thứ nặng: $4 + 3 = 7$ kg.
}
\end{ex}

\begin{ex}
Hai thùng nước mắm có tổng cộng 55 lít, thùng thứ nhất có 30 lít. Thùng thứ hai có bao nhiêu lít nước mắm?
\shortans{25}
\loigiai{
Thùng thứ hai có: $55 - 30 = 25$ lít.
}
\end{ex}

\begin{ex}
Hiệu cân nặng giữa bao gạo 50 kg và bao ngô 35 kg là bao nhiêu ki-lô-gam?
\shortans{15}
\loigiai{
Hiệu: $50 - 35 = 15$ kg.
}
\end{ex}

\begin{ex}
Bác nông dân có 70 lít mật ong, bác đã bán đi 40 lít mật ong. Bác còn lại bao nhiêu lít mật ong?
\shortans{30}
\loigiai{
Còn lại: $70 - 40 = 30$ lít.
}
\end{ex}

\begin{ex}
Một cửa hàng ngày thứ nhất bán được 26 kg đường, ngày thứ hai bán được nhiều hơn ngày thứ nhất 13 kg đường. Ngày thứ hai bán được bao nhiêu ki-lô-gam đường?
\shortans{39}
\loigiai{
Ngày thứ hai bán: $26 + 13 = 39$ kg.
}
\end{ex}

\begin{ex}
Một bình nước có 15 lít nước, người ta rót ra 3 ca, mỗi ca 2 lít. Trong bình còn lại bao nhiêu lít nước?
\shortans{9}
\loigiai{
Rót ra: $2 + 2 + 2 = 6$ lít. Còn lại: $15 - 6 = 9$ lít.
}
\end{ex}

\begin{ex}
Bé Nam nặng 25 kg, khi đeo cặp sách 3 kg lên lưng thì cả người và cặp nặng bao nhiêu ki-lô-gam?
\shortans{28}
\loigiai{
Nặng: $25 + 3 = 28$ kg.
}
\end{ex}

\begin{ex}
Trên cân đĩa, bên trái có 1 bao gạo và 1 quả cân 2 kg, bên phải có 2 quả cân loại 5 kg thì cân thăng bằng. Bao gạo nặng bao nhiêu ki-lô-gam?
\shortans{8}
\loigiai{
Đĩa phải: $5 + 5 = 10$ kg. Bao gạo nặng: $10 - 2 = 8$ kg.
}
\end{ex}

\begin{ex}
Một chiếc thùng rỗng nặng 2 kg, khi đổ vào thùng 15 lít nước (mỗi lít nước nặng 1 kg) thì cả thùng và nước nặng bao nhiêu ki-lô-gam?
\shortans{17}
\loigiai{
Cả thùng nặng: $15 + 2 = 17$ kg.
}
\end{ex}

\begin{ex}
Tính: $10\text{ kg} + 20\text{ kg} + 30\text{ kg} - 15\text{ kg}$.
\shortans{45}
\loigiai{
Ta có $10 + 20 + 30 = 60$, rồi $60 - 15 = 45$ kg.
}
\end{ex}
